\chapter{AI Tool Usage and Development Log}
\label{app:genai-log}

This appendix records the use of an AI coding assistant (Claude Code) during the implementation of the pipeline described in Chapter~\ref{chap:methodology}, in response to the project's own documentation requirement that key development stages record the prompts used, the assistant's responses, and the subsequent modifications and verification steps applied to them. It is a condensed, thesis-appropriate summary. The complete, chronologically ordered development log, every dated entry, including stages not reproduced below, is kept as \texttt{worklog.md} in the project repository, and is the authoritative record if this appendix and the repository log ever appear to diverge.

\section{Scope and General Pattern}
Across the project, the assistant was used only to implement, refactor, or verify code against a specification that had already been decided by hand: a bug taxonomy, a factorial experimental design, a fusion rule and its parameters, a choice of statistical tests, or an explicit refactoring or renaming instruction. It was not used to choose the experimental design, to select which results to report, or to decide the post-hoc reframing of the proposed method described in Section~\ref{sec:results-decisions}. Every substantive use is disclosed below with the prompt that initiated it (verbatim where given in English, translated where the original was in Chinese), a summary of the assistant's contribution, what the author independently verified before relying on the output, and, where relevant, what the assistant proposed that was reviewed and rejected.

\section{Episode 1: Pre-collection pipeline refactor (2026-07-30)}
\textbf{Prompt (translated):} ``before running the formal experiment, open a branch and refactor the code for me, tidy-up only, do not change the method logic; then can you delete the datasets on the remote without touching history?''

\textbf{AI contribution:} Extraction of three shared modules (\texttt{paths.py}, \texttt{conditions.py}, \texttt{stats\_tests.py}) from duplication that had built up across the pipeline's entry-point scripts, and targeted fixes to import-time side effects and dead code in \texttt{build\_manifest.py} and \texttt{score\_trial.py}.

\textbf{Human ownership:} The scope was narrowed, at the point of the request, from ``conservative cleanup only'' to ``structural tidy-up plus extract shared modules,'' a choice made by the author rather than proposed unprompted by the assistant. Running this refactor strictly \emph{before} any EXP002 API call, so no frozen configuration would be edited while results were being collected, followed a pre-existing project rule (Section~\ref{sec:exp_design}) that the request was scheduled to respect, not a judgement made by the assistant.

\textbf{Verification performed:} Every artefact the refactored code could produce was captured before the change and byte-compared after it: the three frozen prompt templates' SHA-256 hashes against the experiment configuration, the manifest files (120 clips), the pre-registered power-analysis output, a full re-score of 1{,}439 synthetic responses against a fixture, dry-run payload files and their decoded-frame \emph{pixel} hashes (so any drift in decoding or frame selection, not just prompt text, would be caught), and the video quality-gate report. All checks matched exactly; four deliberately malformed condition codes, used as negative controls, were confirmed to still be rejected with their original error messages. The dataset-deletion half of the request was scoped down after investigation: the historical dataset blobs are only removable by a git history rewrite, which was identified as out of scope and not performed.

\section{Episode 2: Repository restructuring (2026-07-30)}
\textbf{Prompt (translated):} ``but I think the current file and folder structure is ugly, and it's the formal experiment now but the folders and filenames still say poc, can you deal with that?''

\textbf{AI contribution:} Renaming and relocating directories (\texttt{poc/} to \texttt{pipeline/}, \texttt{latest/} to \texttt{generator/}, hoisting \texttt{data/} and \texttt{results/} to the repository root) using \texttt{git mv} so history was preserved, and amending the path strings inside both experiments' frozen configuration files to match.

\textbf{Human ownership:} The specific choice of new names, and the decision to hoist \texttt{data/} and \texttt{results/} to the top level rather than rename in place, were made explicitly by the author when offered as an alternative. Amending a frozen configuration file is normally prohibited by the project's own freeze policy (Section~\ref{sec:exp_design}); this was judged acceptable here only because EXP002 had issued zero API calls at the time, which was checked before the change was made, and the amendment was recorded with a dated note in the configuration file itself rather than applied silently.

\textbf{Verification performed:} The same golden-output comparison used in Episode 1 was re-run against the pre-refactor baselines in the new directory layout, including the decoded-frame pixel hashes and the prompt SHA-256 values against both experiments' configuration files. One pre-existing mismatch was found and diagnosed, not caused by this change: EXP001's frozen configuration pins an older baseline-prompt hash than the current source. Extracting the literal prompt string from both experiment configurations confirmed this was the deliberate, disclosed prompt revision between EXP001 and EXP002 (Section~\ref{sec:exp_design}), not accidental drift.

\section{Episode 3: EXP002 scoring, reconciliation, and reporting (2026-08-02 to 2026-08-04)}
\textbf{Tool:} Claude Code (Opus 5).

\textbf{Prompts:} An initial request to check a progress email against what the formal trial had actually produced, followed by requests to rewrite the affected passages, to build result tables meeting the project's statistical reporting rule (Section~\ref{sec:benchmarking}), and to judge what belonged in the corresponding commit.

\textbf{AI contribution:} The reconciliation that surfaced three reporting errors before they were sent to the project supervisor (a per-type recall figure attributed to the wrong condition, a false-positive rate quoted against the wrong denominator, and a schema-compliance failure mischaracterised as a parsing failure, all three corrected in Section~\ref{sec:results-discoveries} of this dissertation); the analysis-time discoveries recorded in that same section; and drafting of the corrected passages and dispersion figures subsequently reported.

\textbf{Human ownership:} The decision to reframe the proposed method from A3B1 to A1B1 (Section~\ref{sec:results-decisions}), the decision not to run the additional A1B0 cell, the decision to drop the planned cross-model comparison with Qwen2.5-VL (Section~\ref{sec:env_setup-model}), and the final content of what was reported to the supervisor were all made by the author, not the assistant. One specific assistant proposal, to run A1B0 as a declared supplementary experiment once the fusion-harms-accuracy finding emerged, was considered and rejected, on the grounds given in Section~\ref{sec:results-decisions} that a design change made in direct response to already-observed results needs a new, separately pre-registered experiment rather than an amendment to a frozen one.

\textbf{Verification performed:} Every figure produced in this episode was recomputed independently from the cached raw responses in \texttt{results/exp002/raw/} and checked against the machine-readable scoring summary before being used in any report; no statistic appearing in the project repository or in this dissertation was produced or edited directly by the assistant without this check.

\section{Use in Drafting This Dissertation}
An AI coding assistant was also used in preparing the text of this dissertation, mainly to draft prose expanding Chapter~\ref{chap:methodology}'s methodological detail and Chapter~\ref{chap:results_and_discussion}'s results narrative from the project's own machine-readable outputs (\texttt{results/exp002/summary.json}, the frozen experiment configurations, and \texttt{worklog.md}), and to cross-check the numeric claims made in the text against those sources. Every quantitative claim, table entry, and code-behaviour description in Chapters~\ref{chap:methodology} and~\ref{chap:results_and_discussion} was checked by the author against the corresponding source file before being kept. The interpretive judgements, which findings matter, how the proposed method should be reframed, and what belongs in the limitations, were the author's, following directly from the decisions already recorded in Section~\ref{sec:results-decisions} and Episode 3 above.
